\chapter{Clustering}

\section{Ziel dieses Kapitels}

In den bisherigen Kapiteln haben wir hauptsächlich überwachte Lernverfahren (supervised learning) betrachtet.
Dort gibt es Trainingsdaten mit Eingaben und Zielwerten:

\[
    \mathcal{D}
    =
    \{(\vec{x}_1,t_1),\dots,(\vec{x}_N,t_N)\}.
\]

Beim Clustering ist die Situation anders.
Wir haben nur Eingaben:

\[
    \mathcal{D}
    =
    \{\vec{x}_1,\dots,\vec{x}_N\}.
\]

Es gibt keine beobachteten Zielwerte \(t_n\).
Das Ziel ist, Struktur in den Daten zu finden.

\begin{conceptbox}
Clustering ist ein unüberwachtes Lernverfahren (unsupervised learning).

Das Ziel ist, Datenpunkte in Gruppen einzuteilen, sodass Punkte innerhalb einer Gruppe ähnlich sind und Punkte aus verschiedenen Gruppen möglichst unterschiedlich sind.
\end{conceptbox}

Nach diesem Kapitel solltest du folgende Fragen beantworten können:

\begin{itemize}
    \item Was ist Clustering?
    \item Was ist der Unterschied zwischen Klassifikation und Clustering?
    \item Was bedeutet Ähnlichkeit oder Distanz zwischen Datenpunkten?
    \item Was ist \(k\)-means?
    \item Wie funktioniert der \(k\)-means-Algorithmus?
    \item Welche Zielfunktion minimiert \(k\)-means?
    \item Warum besteht \(k\)-means aus einem Zuordnungsschritt und einem Aktualisierungsschritt?
    \item Warum konvergiert \(k\)-means?
    \item Warum hängt \(k\)-means von der Initialisierung ab?
    \item Wie wählt man die Anzahl der Cluster \(K\)?
    \item Was sind Elbow-Methode und Silhouette-Koeffizient?
    \item Was sind hierarchisches Clustering und DBSCAN?
    \item Welche Grenzen hat Clustering?
\end{itemize}

\section{Was ist Clustering?}

Clustering bezeichnet Methoden, die Datenpunkte automatisch in Gruppen einteilen.
Diese Gruppen nennt man Cluster.

Gegeben seien Datenpunkte

\[
    \vec{x}_1,\dots,\vec{x}_N
    \in
    \mathbb{R}^D.
\]

Ein Clustering ordnet jedem Datenpunkt eine Clusterzugehörigkeit zu.
Wenn es \(K\) Cluster gibt, dann ist die Zuordnung:

\[
    z_n \in \{1,\dots,K\}.
\]

Dabei bedeutet

\[
    z_n=k,
\]

dass der Datenpunkt \(\vec{x}_n\) zum Cluster \(k\) gehört.

\begin{definitionbox}
Ein Clustering ist eine Zuordnung

\[
    z_n \in \{1,\dots,K\}
\]

für jeden Datenpunkt \(\vec{x}_n\).

Der Wert \(z_n\) gibt an, zu welchem Cluster der Datenpunkt gehört.
\end{definitionbox}

\begin{examplebox}
Beispiele für Clustering:

\begin{itemize}
    \item Kundinnen und Kunden in Gruppen ähnlichen Kaufverhaltens einteilen.
    \item Patientinnen und Patienten mit ähnlichen Messwerten gruppieren.
    \item Bilder nach visueller Ähnlichkeit gruppieren.
    \item Dokumente nach Themen gruppieren.
    \item Messdaten in typische Zustände eines Systems aufteilen.
\end{itemize}
\end{examplebox}

\section{Clustering versus Klassifikation}

Clustering und Klassifikation wirken ähnlich, sind aber grundlegend verschieden.

Bei Klassifikation (classification) gibt es Labels:

\[
    t_n \in \{1,\dots,K\}.
\]

Das Modell lernt aus Beispielen mit bekannten Klassen.

Bei Clustering gibt es keine Labels.
Das Modell muss Gruppen selbst finden.

\begin{definitionbox}
Klassifikation ist überwachtes Lernen (supervised learning):

\[
    (\vec{x}_n,t_n)
\]

sind gegeben.

Clustering ist unüberwachtes Lernen (unsupervised learning):

\[
    \vec{x}_n
\]

sind gegeben, aber keine Labels \(t_n\).
\end{definitionbox}

\begin{conceptbox}
Klassifikation beantwortet:

\[
    \text{Welche bekannte Klasse hat dieser Datenpunkt?}
\]

Clustering beantwortet:

\[
    \text{Welche Gruppenstruktur scheint in den Daten vorhanden zu sein?}
\]
\end{conceptbox}

\section{Ähnlichkeit und Distanz}

Clustering braucht eine Vorstellung davon, wann Datenpunkte ähnlich sind.
Oft wird Ähnlichkeit über eine Distanz definiert.

Für zwei Datenpunkte

\[
    \vec{x}_n,\vec{x}_m\in\mathbb{R}^D
\]

ist die euklidische Distanz:

\[
    d(\vec{x}_n,\vec{x}_m)
    =
    \|\vec{x}_n-\vec{x}_m\|.
\]

Ausgeschrieben:

\[
    \|\vec{x}_n-\vec{x}_m\|
    =
    \sqrt{
        \sum_{j=1}^{D}
        (x_{nj}-x_{mj})^2
    }.
\]

Häufig verwendet man auch die quadrierte euklidische Distanz:

\[
    \|\vec{x}_n-\vec{x}_m\|^2
    =
    \sum_{j=1}^{D}
    (x_{nj}-x_{mj})^2.
\]

\begin{definitionbox}
Die euklidische Distanz zwischen zwei Datenpunkten ist:

\[
    d(\vec{x}_n,\vec{x}_m)
    =
    \sqrt{
        \sum_{j=1}^{D}
        (x_{nj}-x_{mj})^2
    }.
\]
\end{definitionbox}

\begin{notebox}
Welche Distanz sinnvoll ist, hängt vom Problem ab.

Für numerische Vektoren ist die euklidische Distanz häufig.
Für Texte, Sequenzen oder kategoriale Daten können andere Ähnlichkeitsmaße sinnvoller sein.
\end{notebox}

\section{Feature-Skalierung}

Clustering ist oft stark abhängig von der Skalierung der Merkmale.
Wenn ein Merkmal sehr große Zahlenwerte hat, dominiert es die Distanz.

\begin{examplebox}
Angenommen, ein Datenpunkt enthält:

\[
    x_1 = \text{Alter in Jahren}
\]

und

\[
    x_2 = \text{Einkommen in Euro}.
\]

Dann kann das Einkommen numerisch viel größere Werte haben als das Alter.
Ohne Skalierung wird die euklidische Distanz fast nur durch das Einkommen bestimmt.
\end{examplebox}

Eine Standardisierung ist:

\[
    x_j'
    =
    \frac{x_j-\mu_j}{\sigma_j}.
\]

Dabei sind \(\mu_j\) und \(\sigma_j\) Mittelwert und Standardabweichung des Merkmals \(j\).

\begin{examtipbox}
Vor distanzbasiertem Clustering sollte man prüfen, ob die Merkmale sinnvoll skaliert sind.

Bei unterschiedlichen Einheiten oder Größenordnungen ist Standardisierung oft notwendig.
\end{examtipbox}

\section{\(k\)-means: Grundidee}

Der bekannteste Clustering-Algorithmus ist \(k\)-means.
Dabei wird die Anzahl der Cluster \(K\) vorgegeben.

Die Idee ist:
Jeder Cluster wird durch einen Mittelpunkt beschrieben.
Diese Mittelpunkte heißen Clusterzentren (cluster centers) oder Zentroiden (centroids).

Wir schreiben die Zentren als

\[
    \vec{\mu}_1,\dots,\vec{\mu}_K
    \in
    \mathbb{R}^D.
\]

Jeder Datenpunkt wird dem nächstgelegenen Zentrum zugeordnet.

\begin{definitionbox}
Beim \(k\)-means-Clustering gibt es \(K\) Clusterzentren

\[
    \vec{\mu}_1,\dots,\vec{\mu}_K.
\]

Jeder Datenpunkt \(\vec{x}_n\) wird dem nächstgelegenen Zentrum zugeordnet.
\end{definitionbox}

\begin{conceptbox}
\(k\)-means versucht, Cluster so zu finden, dass Datenpunkte möglichst nahe am Zentrum ihres Clusters liegen.
\end{conceptbox}

\section{Cluster-Zuordnungen}

Für jeden Datenpunkt \(\vec{x}_n\) gibt es eine Zuordnung

\[
    z_n\in\{1,\dots,K\}.
\]

Wenn

\[
    z_n=k,
\]

dann gehört \(\vec{x}_n\) zu Cluster \(k\).

Die Zuordnung bei gegebenen Zentren ist:

\[
    z_n
    =
    \arg\min_{k\in\{1,\dots,K\}}
    \|\vec{x}_n-\vec{\mu}_k\|^2.
\]

\begin{definitionbox}
Der Zuordnungsschritt bei \(k\)-means ist:

\[
    z_n
    =
    \arg\min_k
    \|\vec{x}_n-\vec{\mu}_k\|^2.
\]
\end{definitionbox}

Man kann statt der quadrierten Distanz auch die normale euklidische Distanz verwenden.
Da die Wurzelfunktion monoton wachsend ist, ist das nächstgelegene Zentrum dasselbe:

\[
    \arg\min_k
    \|\vec{x}_n-\vec{\mu}_k\|
    =
    \arg\min_k
    \|\vec{x}_n-\vec{\mu}_k\|^2.
\]

\section{Zielfunktion von \(k\)-means}

\(k\)-means minimiert die Summe der quadrierten Abstände jedes Datenpunkts zu seinem Clusterzentrum.

\[
    J
    =
    \sum_{n=1}^{N}
    \|\vec{x}_n-\vec{\mu}_{z_n}\|^2.
\]

Diese Größe heißt innerhalb-Cluster-Quadratsumme (within-cluster sum of squares).

\begin{definitionbox}
Die \(k\)-means-Zielfunktion ist:

\[
    J(\{z_n\},\{\vec{\mu}_k\})
    =
    \sum_{n=1}^{N}
    \|\vec{x}_n-\vec{\mu}_{z_n}\|^2.
\]

Dabei ist \(z_n\) die Clusterzuordnung von \(\vec{x}_n\).
\end{definitionbox}

Äquivalent kann man mit Indikatorvariablen schreiben:

\[
    r_{nk}
    =
    \begin{cases}
        1, & z_n=k, \\
        0, & z_n\neq k.
    \end{cases}
\]

Dann ist:

\[
    J
    =
    \sum_{n=1}^{N}
    \sum_{k=1}^{K}
    r_{nk}
    \|\vec{x}_n-\vec{\mu}_k\|^2.
\]

\begin{notebox}
Für jeden Datenpunkt gilt:

\[
    \sum_{k=1}^{K} r_{nk}=1.
\]

Jeder Punkt gehört also genau zu einem Cluster.
\end{notebox}

\section{Der \(k\)-means-Algorithmus}

Die Zielfunktion hängt von zwei Arten von Variablen ab:

\begin{itemize}
    \item Zuordnungen \(z_n\) oder \(r_{nk}\),
    \item Clusterzentren \(\vec{\mu}_k\).
\end{itemize}

Wenn die Zentren fest sind, sind die optimalen Zuordnungen leicht zu bestimmen.
Wenn die Zuordnungen fest sind, sind die optimalen Zentren leicht zu bestimmen.

Der \(k\)-means-Algorithmus wechselt daher zwischen zwei Schritten:

\begin{enumerate}
    \item Zuordnungsschritt (assignment step),
    \item Aktualisierungsschritt (update step).
\end{enumerate}

\begin{definitionbox}
Der \(k\)-means-Algorithmus:

\begin{enumerate}
    \item Wähle initiale Zentren \(\vec{\mu}_1,\dots,\vec{\mu}_K\).
    \item Ordne jeden Datenpunkt dem nächsten Zentrum zu:
    \[
        z_n
        =
        \arg\min_k
        \|\vec{x}_n-\vec{\mu}_k\|^2.
    \]
    \item Aktualisiere jedes Zentrum als Mittelwert der zugeordneten Punkte:
    \[
        \vec{\mu}_k
        =
        \frac{1}{N_k}
        \sum_{n:z_n=k}
        \vec{x}_n.
    \]
    \item Wiederhole Schritt \(2\) und \(3\), bis sich die Zuordnungen oder Zentren nicht mehr ändern.
\end{enumerate}
\end{definitionbox}

Dabei ist

\[
    N_k
    =
    \sum_{n=1}^{N}
    r_{nk}
\]

die Anzahl der Punkte in Cluster \(k\).

\section{Herleitung des Aktualisierungsschritts}

Wir nehmen an, dass die Zuordnungen \(r_{nk}\) fest sind.
Dann minimieren wir

\[
    J
    =
    \sum_{n=1}^{N}
    \sum_{k=1}^{K}
    r_{nk}
    \|\vec{x}_n-\vec{\mu}_k\|^2
\]

nach den Zentren \(\vec{\mu}_k\).

Für ein bestimmtes Cluster \(k\) betrachten wir nur die Terme, die \(\vec{\mu}_k\) enthalten:

\[
    J_k
    =
    \sum_{n=1}^{N}
    r_{nk}
    \|\vec{x}_n-\vec{\mu}_k\|^2.
\]

Wir leiten nach \(\vec{\mu}_k\) ab:

\[
    \frac{\partial J_k}{\partial \vec{\mu}_k}
    =
    \sum_{n=1}^{N}
    r_{nk}
    2(\vec{\mu}_k-\vec{x}_n).
\]

Setze gleich \(\vec{0}\):

\[
    \sum_{n=1}^{N}
    r_{nk}
    2(\vec{\mu}_k-\vec{x}_n)
    =
    \vec{0}.
\]

Teilen durch \(2\):

\[
    \sum_{n=1}^{N}
    r_{nk}
    (\vec{\mu}_k-\vec{x}_n)
    =
    \vec{0}.
\]

Ausmultiplizieren:

\[
    \left(
        \sum_{n=1}^{N}
        r_{nk}
    \right)
    \vec{\mu}_k
    -
    \sum_{n=1}^{N}
    r_{nk}\vec{x}_n
    =
    \vec{0}.
\]

Also:

\[
    \left(
        \sum_{n=1}^{N}
        r_{nk}
    \right)
    \vec{\mu}_k
    =
    \sum_{n=1}^{N}
    r_{nk}\vec{x}_n.
\]

Damit:

\[
    \vec{\mu}_k
    =
    \frac{
        \sum_{n=1}^{N}
        r_{nk}\vec{x}_n
    }{
        \sum_{n=1}^{N}
        r_{nk}
    }.
\]

Da

\[
    N_k
    =
    \sum_{n=1}^{N}
    r_{nk},
\]

ist:

\[
    \vec{\mu}_k
    =
    \frac{1}{N_k}
    \sum_{n:z_n=k}
    \vec{x}_n.
\]

\begin{conceptbox}
Das optimale Zentrum eines Clusters ist der Mittelwert aller Datenpunkte in diesem Cluster.
\end{conceptbox}

\section{Warum \(k\)-means konvergiert}

In jedem Zuordnungsschritt wird jeder Punkt dem nächstgelegenen Zentrum zugeordnet.
Dadurch kann die Zielfunktion \(J\) nicht steigen.

In jedem Aktualisierungsschritt wird jedes Zentrum als Mittelwert seiner Punkte gesetzt.
Das minimiert \(J\) für feste Zuordnungen.
Auch dadurch kann \(J\) nicht steigen.

Also gilt nach jeder Iteration:

\[
    J^{(\tau+1)}
    \leq
    J^{(\tau)}.
\]

Da

\[
    J \geq 0,
\]

kann \(J\) nicht unbegrenzt fallen.

\begin{conceptbox}
\(k\)-means konvergiert, weil die Zielfunktion in jedem Schritt nicht steigt und nach unten durch \(0\) beschränkt ist.
\end{conceptbox}

\begin{notebox}
Konvergenz bedeutet hier:
Der Algorithmus erreicht ein lokales Minimum oder einen festen Punkt.

Es bedeutet nicht, dass immer das globale Minimum gefunden wird.
\end{notebox}

\section{Lokale Minima und Initialisierung}

Die \(k\)-means-Zielfunktion ist nicht konvex in allen Variablen gemeinsam.
Der Algorithmus kann daher in verschiedenen lokalen Minima landen.

Das Ergebnis hängt stark von der Initialisierung der Zentren ab.

\begin{examplebox}
Wenn zwei initiale Zentren zufällig sehr nahe beieinander liegen, können sie anfangs denselben Bereich der Daten erklären.
Ein anderer Bereich bekommt dann vielleicht kein gutes Zentrum.
Das kann zu einem schlechteren lokalen Minimum führen.
\end{examplebox}

\begin{conceptbox}
\(k\)-means sollte oft mehrfach mit verschiedenen Initialisierungen gestartet werden.

Dann wählt man die Lösung mit dem kleinsten Wert der Zielfunktion \(J\).
\end{conceptbox}

\section{\(k\)-means++}

Eine bessere Initialisierungsmethode ist \(k\)-means++.

Die Idee:
Zentren sollen nicht nur zufällig gewählt werden, sondern möglichst gut über den Datensatz verteilt sein.

Vereinfacht:

\begin{enumerate}
    \item Wähle das erste Zentrum zufällig aus den Datenpunkten.
    \item Für jeden Datenpunkt berechne die Distanz zum nächsten bereits gewählten Zentrum.
    \item Wähle das nächste Zentrum mit Wahrscheinlichkeit proportional zur quadrierten Distanz.
    \item Wiederhole, bis \(K\) Zentren gewählt sind.
\end{enumerate}

\begin{conceptbox}
\(k\)-means++ wählt initiale Zentren tendenziell weit voneinander entfernt.

Dadurch startet \(k\)-means oft in einer besseren Ausgangslage.
\end{conceptbox}

\section{Leere Cluster}

Während des \(k\)-means-Algorithmus kann es passieren, dass einem Cluster kein Datenpunkt zugeordnet wird.
Dann ist

\[
    N_k=0.
\]

Die Aktualisierung

\[
    \vec{\mu}_k
    =
    \frac{1}{N_k}
    \sum_{n:z_n=k}
    \vec{x}_n
\]

ist dann nicht definiert.

\begin{notebox}
Leere Cluster müssen speziell behandelt werden.

Mögliche Strategien:

\begin{itemize}
    \item Zentrum zufällig neu initialisieren.
    \item Zentrum auf einen weit entfernten Datenpunkt setzen.
    \item Cluster entfernen, falls weniger als \(K\) Cluster erlaubt sind.
\end{itemize}
\end{notebox}

\section{Wahl der Clusterzahl \(K\)}

Bei \(k\)-means muss \(K\) vorgegeben werden.
Das ist oft schwierig.

Wenn \(K\) zu klein ist, werden unterschiedliche Gruppen zusammengelegt.
Wenn \(K\) zu groß ist, werden natürliche Gruppen künstlich aufgeteilt.

\begin{conceptbox}
Die Wahl von \(K\) ist ein Modellselektionsproblem.

Es gibt selten eine eindeutig richtige Clusterzahl.
Die sinnvolle Wahl hängt von Daten, Ziel und Interpretation ab.
\end{conceptbox}

\section{Elbow-Methode}

Die Elbow-Methode betrachtet die \(k\)-means-Zielfunktion als Funktion von \(K\).

Für jedes \(K\) berechnet man:

\[
    J_K
    =
    \sum_{n=1}^{N}
    \|\vec{x}_n-\vec{\mu}_{z_n}\|^2.
\]

Wenn \(K\) größer wird, kann \(J_K\) nur sinken oder gleich bleiben.
Mehr Cluster können die Daten immer mindestens genauso gut erklären.

Man zeichnet \(J_K\) gegen \(K\).
Oft sucht man nach einem Knickpunkt, dem sogenannten Ellenbogen.

\begin{definitionbox}
Die Elbow-Methode wählt \(K\) anhand eines Knickpunkts in der Kurve

\[
    K
    \longmapsto
    J_K.
\]
\end{definitionbox}

\begin{notebox}
Die Elbow-Methode ist heuristisch.

Manchmal gibt es keinen klaren Knick.
Dann ist die Wahl von \(K\) nicht eindeutig.
\end{notebox}

\section{Silhouette-Koeffizient}

Der Silhouette-Koeffizient misst, wie gut ein Datenpunkt zu seinem eigenen Cluster passt im Vergleich zu anderen Clustern.

Für einen Datenpunkt \(\vec{x}_n\):

\begin{itemize}
    \item \(a_n\): durchschnittliche Distanz zu Punkten im eigenen Cluster,
    \item \(b_n\): kleinste durchschnittliche Distanz zu Punkten in einem anderen Cluster.
\end{itemize}

Der Silhouette-Wert ist:

\[
    s_n
    =
    \frac{b_n-a_n}{\max(a_n,b_n)}.
\]

Es gilt:

\[
    -1 \leq s_n \leq 1.
\]

\begin{definitionbox}
Der Silhouette-Koeffizient eines Datenpunkts ist:

\[
    s_n
    =
    \frac{b_n-a_n}{\max(a_n,b_n)}.
\]

Dabei misst \(a_n\) die Nähe zum eigenen Cluster und \(b_n\) die Nähe zum nächsten anderen Cluster.
\end{definitionbox}

Interpretation:

\begin{itemize}
    \item \(s_n\approx 1\): Punkt passt gut zum eigenen Cluster.
    \item \(s_n\approx 0\): Punkt liegt zwischen Clustern.
    \item \(s_n<0\): Punkt passt besser zu einem anderen Cluster.
\end{itemize}

Der mittlere Silhouette-Wert ist:

\[
    s
    =
    \frac{1}{N}
    \sum_{n=1}^{N}
    s_n.
\]

\begin{notebox}
Der Silhouette-Koeffizient kann helfen, \(K\) zu wählen.
Aber auch er ist nur ein internes Qualitätsmaß und ersetzt keine fachliche Interpretation.
\end{notebox}

\section{Grenzen von \(k\)-means}

\(k\)-means ist einfach und effizient, hat aber wichtige Einschränkungen.

\subsection{Sphärische Cluster}

\(k\)-means funktioniert besonders gut, wenn Cluster ungefähr kugelförmig sind und ähnliche Größe haben.

Das liegt an der euklidischen Distanz zu Zentren.

\subsection{Empfindlichkeit gegenüber Ausreißern}

Da Mittelwerte verwendet werden, können Ausreißer die Zentren stark verschieben.

\subsection{Feste Clusterzahl}

Die Anzahl \(K\) muss vorher festgelegt werden.

\subsection{Harte Zuordnungen}

Jeder Punkt gehört genau zu einem Cluster.
Unsicherheit wird nicht modelliert.

\subsection{Skalierung}

Unterschiedliche Merkmalsskalen können das Ergebnis stark verzerren.

\begin{examtipbox}
Typische Grenzen von \(k\)-means:

\begin{itemize}
    \item \(K\) muss vorgegeben werden.
    \item Ergebnis hängt von Initialisierung ab.
    \item Cluster sollten ungefähr konvex und kugelförmig sein.
    \item Ausreißer können starke Effekte haben.
    \item Feature-Skalierung ist wichtig.
    \item Es gibt keine probabilistische Unsicherheit der Zuordnung.
\end{itemize}
\end{examtipbox}

\section{Harte und weiche Clusterzuordnung}

Bei \(k\)-means ist die Zuordnung hart:

\[
    r_{nk}\in\{0,1\}.
\]

Ein Datenpunkt gehört genau zu einem Cluster.

Bei weichem Clustering (soft clustering) kann ein Punkt teilweise mehreren Clustern zugeordnet sein.
Dann sind die Verantwortlichkeiten (responsibilities):

\[
    r_{nk}\in[0,1],
\]

mit

\[
    \sum_{k=1}^{K}
    r_{nk}
    =
    1.
\]

\begin{definitionbox}
Hartes Clustering (hard clustering) ordnet jeden Punkt genau einem Cluster zu.

Weiches Clustering (soft clustering) ordnet jedem Punkt Wahrscheinlichkeiten oder Verantwortlichkeiten für mehrere Cluster zu.
\end{definitionbox}

\begin{conceptbox}
Gaussian Mixture Models verwenden weiche Clusterzuordnungen.

\(k\)-means kann als harte und vereinfachte Variante einer Mischungsidee verstanden werden.
\end{conceptbox}

\section{Hierarchisches Clustering}

Hierarchisches Clustering erzeugt nicht nur eine einzelne Clusterlösung, sondern eine Hierarchie von Clustern.

Es gibt zwei Haupttypen:

\begin{itemize}
    \item agglomeratives Clustering (agglomerative clustering),
    \item divisives Clustering (divisive clustering).
\end{itemize}

Beim agglomerativen Clustering startet jeder Datenpunkt als eigener Cluster.
Dann werden schrittweise die ähnlichsten Cluster zusammengeführt.

Beim divisiven Clustering startet man mit einem großen Cluster und teilt ihn schrittweise auf.

\begin{definitionbox}
Hierarchisches Clustering erzeugt eine verschachtelte Struktur von Clustern.

Diese Struktur kann als Dendrogramm dargestellt werden.
\end{definitionbox}

\section{Agglomeratives Clustering}

Agglomeratives Clustering läuft so ab:

\begin{enumerate}
    \item Jeder Datenpunkt ist zunächst ein eigener Cluster.
    \item Berechne Distanzen zwischen Clustern.
    \item Führe die zwei ähnlichsten Cluster zusammen.
    \item Aktualisiere die Distanzen.
    \item Wiederhole, bis nur noch ein Cluster übrig ist.
\end{enumerate}

\begin{conceptbox}
Agglomeratives Clustering baut eine Clusterhierarchie von unten nach oben auf.

\[
    N \text{ einzelne Cluster}
    \quad
    \longrightarrow
    \quad
    1 \text{ großer Cluster}.
\]
\end{conceptbox}

\section{Linkage-Kriterien}

Beim hierarchischen Clustering muss definiert werden, wie die Distanz zwischen zwei Clustern gemessen wird.
Das nennt man Linkage-Kriterium.

Seien \(A\) und \(B\) zwei Cluster.

\subsection{Single Linkage}

\[
    d_{\mathrm{single}}(A,B)
    =
    \min_{\vec{x}\in A,\vec{x}'\in B}
    d(\vec{x},\vec{x}').
\]

Single Linkage verwendet den kleinsten Abstand zwischen Punkten der beiden Cluster.

\subsection{Complete Linkage}

\[
    d_{\mathrm{complete}}(A,B)
    =
    \max_{\vec{x}\in A,\vec{x}'\in B}
    d(\vec{x},\vec{x}').
\]

Complete Linkage verwendet den größten Abstand.

\subsection{Average Linkage}

\[
    d_{\mathrm{average}}(A,B)
    =
    \frac{1}{|A||B|}
    \sum_{\vec{x}\in A}
    \sum_{\vec{x}'\in B}
    d(\vec{x},\vec{x}').
\]

Average Linkage verwendet den durchschnittlichen Abstand.

\subsection{Ward-Linkage}

Ward-Linkage führt diejenigen Cluster zusammen, bei denen die Zunahme der innerhalb-Cluster-Quadratsumme am kleinsten ist.

\begin{definitionbox}
Ein Linkage-Kriterium definiert die Distanz zwischen Clustern.

Es bestimmt, welche Cluster beim hierarchischen Clustering zusammengeführt werden.
\end{definitionbox}

\begin{notebox}
Unterschiedliche Linkage-Kriterien können sehr unterschiedliche Clusterstrukturen ergeben.
\end{notebox}

\section{Dendrogramm}

Ein Dendrogramm ist eine Baumdarstellung der Clusterhierarchie.
Die Blätter sind einzelne Datenpunkte.
Innere Knoten entsprechen Zusammenführungen von Clustern.

Die Höhe einer Zusammenführung zeigt typischerweise die Distanz, bei der die Cluster zusammengeführt wurden.

\begin{conceptbox}
Ein Dendrogramm erlaubt, die Anzahl der Cluster nachträglich zu wählen.

Man schneidet den Baum auf einer bestimmten Höhe.
\end{conceptbox}

Wenn man das Dendrogramm hoch schneidet, erhält man wenige große Cluster.
Wenn man es niedrig schneidet, erhält man viele kleine Cluster.

\section{DBSCAN}

DBSCAN steht für Density-Based Spatial Clustering of Applications with Noise.
Es ist ein dichtebasiertes Clustering-Verfahren.

Die Idee:
Cluster sind Bereiche hoher Punktdichte, getrennt durch Bereiche niedriger Punktdichte.

DBSCAN benötigt zwei wichtige Hyperparameter:

\[
    \varepsilon
\]

und

\[
    \mathrm{MinPts}.
\]

\begin{definitionbox}
DBSCAN definiert Cluster als zusammenhängende Regionen hoher Dichte.

Punkte in dünn besetzten Regionen können als Rauschen oder Ausreißer markiert werden.
\end{definitionbox}

\section{Kernpunkte, Randpunkte und Rauschen}

Für einen Punkt \(\vec{x}_n\) betrachtet man seine \(\varepsilon\)-Nachbarschaft:

\[
    \mathcal{N}_{\varepsilon}(\vec{x}_n)
    =
    \{\vec{x}_m : d(\vec{x}_n,\vec{x}_m)\leq \varepsilon\}.
\]

Ein Kernpunkt (core point) ist ein Punkt, dessen \(\varepsilon\)-Nachbarschaft mindestens \(\mathrm{MinPts}\) Punkte enthält.

\[
    |\mathcal{N}_{\varepsilon}(\vec{x}_n)|
    \geq
    \mathrm{MinPts}.
\]

Ein Randpunkt (border point) ist kein Kernpunkt, liegt aber in der Nachbarschaft eines Kernpunkts.

Ein Rauschpunkt (noise point) ist weder Kernpunkt noch Randpunkt.

\begin{definitionbox}
DBSCAN unterscheidet:

\begin{itemize}
    \item Kernpunkte (core points),
    \item Randpunkte (border points),
    \item Rauschpunkte (noise points).
\end{itemize}
\end{definitionbox}

\section{Eigenschaften von DBSCAN}

DBSCAN hat andere Eigenschaften als \(k\)-means.

\begin{itemize}
    \item Die Anzahl der Cluster muss nicht vorgegeben werden.
    \item Cluster können nicht-kugelförmig sein.
    \item Ausreißer können als Rauschen erkannt werden.
    \item Die Wahl von \(\varepsilon\) und \(\mathrm{MinPts}\) ist wichtig.
    \item DBSCAN hat Probleme, wenn Cluster sehr unterschiedliche Dichten haben.
\end{itemize}

\begin{conceptbox}
DBSCAN ist besonders nützlich, wenn Cluster unregelmäßige Formen haben und Ausreißer erkannt werden sollen.
\end{conceptbox}

\section{Vergleich wichtiger Clustering-Methoden}

\[
\begin{array}{c|c|c|c}
\text{Methode}
&
\text{Clusterform}
&
\text{Braucht } K?
&
\text{Ausreißer}
\\
\hline
k\text{-means}
&
\text{kugelförmig}
&
\text{ja}
&
\text{empfindlich}
\\
\hline
\text{Hierarchisch}
&
\text{abhängig von Linkage}
&
\text{nicht direkt}
&
\text{abhängig von Methode}
\\
\hline
\text{DBSCAN}
&
\text{beliebiger}
&
\text{nein}
&
\text{kann Rauschen erkennen}
\end{array}
\]

\begin{notebox}
Es gibt nicht den einen besten Clustering-Algorithmus.

Die passende Methode hängt von Datenstruktur, Ziel und gewünschter Interpretation ab.
\end{notebox}

\section{Clustering und probabilistische Modelle}

\(k\)-means ordnet Punkte hart zu Clustern zu.
Eine probabilistische Alternative sind Mischmodelle (mixture models).

Ein Gaussian Mixture Model nimmt an, dass Daten aus einer Mischung von Normalverteilungen entstehen:

\[
    p(\vec{x})
    =
    \sum_{k=1}^{K}
    \pi_k
    \mathcal{N}(\vec{x}\mid\vec{\mu}_k,\Sigma_k).
\]

Dabei ist:

\begin{itemize}
    \item \(\pi_k\) das Mischgewicht von Komponente \(k\),
    \item \(\vec{\mu}_k\) der Mittelwert,
    \item \(\Sigma_k\) die Kovarianzmatrix.
\end{itemize}

\begin{conceptbox}
Gaussian Mixture Models verallgemeinern \(k\)-means:

\begin{itemize}
    \item weiche Zuordnungen statt harte Zuordnungen,
    \item ellipsenförmige Cluster statt nur kugelförmige Cluster,
    \item probabilistische Interpretation.
\end{itemize}
\end{conceptbox}

Gaussian Mixture Models werden typischerweise mit dem EM-Algorithmus trainiert.
Das wird in den nächsten Kapiteln behandelt.

\section{Evaluation von Clustering}

Clustering ist schwer zu evaluieren, weil keine Labels vorhanden sind.
Man unterscheidet interne und externe Maße.

\subsection{Interne Maße}

Interne Maße verwenden nur die Daten und das Clustering.
Beispiele:

\begin{itemize}
    \item \(k\)-means-Zielfunktion \(J\),
    \item Silhouette-Koeffizient,
    \item innerhalb-Cluster-Distanzen,
    \item zwischen-Cluster-Distanzen.
\end{itemize}

\subsection{Externe Maße}

Externe Maße verwenden bekannte Labels, falls sie für Evaluation verfügbar sind.
Dann kann man Clustering mit echten Klassen vergleichen.

Beispiele:

\begin{itemize}
    \item Adjusted Rand Index,
    \item Normalized Mutual Information,
    \item Purity.
\end{itemize}

\begin{notebox}
Wenn echte Labels vorhanden sind, ist die Aufgabe nicht mehr rein unüberwacht.
Man kann Labels aber verwenden, um ein Clustering nachträglich zu bewerten.
\end{notebox}

\section{Interpretation von Clustern}

Ein Clustering-Algorithmus findet immer irgendeine Struktur.
Das bedeutet nicht automatisch, dass diese Struktur fachlich sinnvoll ist.

\begin{conceptbox}
Cluster müssen interpretiert werden.

Ein mathematisch gutes Clustering ist nicht automatisch ein fachlich sinnvolles Clustering.
\end{conceptbox}

Wichtige Fragen:

\begin{itemize}
    \item Sind die Cluster stabil gegenüber kleinen Änderungen der Daten?
    \item Haben die Cluster eine fachliche Bedeutung?
    \item Entstehen die Cluster nur durch Skalierung oder Vorverarbeitung?
    \item Sind Ausreißer verantwortlich für die Clusterstruktur?
    \item Ist die gewählte Distanz sinnvoll?
\end{itemize}

\section{Clustering-Workflow}

Ein typischer Clustering-Workflow ist:

\begin{enumerate}
    \item Daten verstehen.
    \item Fehlende Werte behandeln.
    \item Merkmale sinnvoll auswählen.
    \item Merkmale skalieren oder standardisieren.
    \item Distanzmaß wählen.
    \item Clustering-Methode wählen.
    \item Hyperparameter wählen.
    \item Clustering durchführen.
    \item Cluster visualisieren.
    \item Cluster fachlich interpretieren.
    \item Stabilität prüfen.
\end{enumerate}

\begin{examtipbox}
Bei Clustering ist Vorverarbeitung besonders wichtig.

Da es keine Labels gibt, kann ein Algorithmus nicht wissen, welche Merkmale fachlich relevant sind.
\end{examtipbox}

\section{Mini-Aufgaben}

\subsection{Aufgabe 1: Nächstes Zentrum}

Gegeben seien zwei Clusterzentren:

\[
    \vec{\mu}_1
    =
    \begin{pmatrix}
        0 \\
        0
    \end{pmatrix},
    \qquad
    \vec{\mu}_2
    =
    \begin{pmatrix}
        4 \\
        0
    \end{pmatrix}.
\]

Ein Datenpunkt sei:

\[
    \vec{x}
    =
    \begin{pmatrix}
        1 \\
        0
    \end{pmatrix}.
\]

Welchem Zentrum wird \(\vec{x}\) bei \(k\)-means zugeordnet?

\begin{examplebox}
Distanz zu \(\vec{\mu}_1\):

\[
    \|\vec{x}-\vec{\mu}_1\|^2
    =
    \left\|
    \begin{pmatrix}
        1 \\
        0
    \end{pmatrix}
    -
    \begin{pmatrix}
        0 \\
        0
    \end{pmatrix}
    \right\|^2
    =
    1^2+0^2
    =
    1.
\]

Distanz zu \(\vec{\mu}_2\):

\[
    \|\vec{x}-\vec{\mu}_2\|^2
    =
    \left\|
    \begin{pmatrix}
        1 \\
        0
    \end{pmatrix}
    -
    \begin{pmatrix}
        4 \\
        0
    \end{pmatrix}
    \right\|^2
    =
    (-3)^2+0^2
    =
    9.
\]

Da

\[
    1<9,
\]

wird \(\vec{x}\) Cluster \(1\) zugeordnet.
\end{examplebox}

\subsection{Aufgabe 2: Zentrum aktualisieren}

Ein Cluster enthält die Punkte:

\[
    \vec{x}_1
    =
    \begin{pmatrix}
        1 \\
        2
    \end{pmatrix},
    \qquad
    \vec{x}_2
    =
    \begin{pmatrix}
        3 \\
        4
    \end{pmatrix}.
\]

Berechne das neue \(k\)-means-Zentrum.

\begin{examplebox}
Das Zentrum ist der Mittelwert:

\[
    \vec{\mu}
    =
    \frac{1}{2}
    \left(
        \vec{x}_1+\vec{x}_2
    \right).
\]

Also:

\[
    \vec{\mu}
    =
    \frac{1}{2}
    \left[
    \begin{pmatrix}
        1 \\
        2
    \end{pmatrix}
    +
    \begin{pmatrix}
        3 \\
        4
    \end{pmatrix}
    \right].
\]

Damit:

\[
    \vec{\mu}
    =
    \frac{1}{2}
    \begin{pmatrix}
        4 \\
        6
    \end{pmatrix}
    =
    \begin{pmatrix}
        2 \\
        3
    \end{pmatrix}.
\]
\end{examplebox}

\subsection{Aufgabe 3: \(k\)-means-Zielfunktion}

Gegeben seien zwei Punkte und ihr Clusterzentrum:

\[
    \vec{x}_1
    =
    \begin{pmatrix}
        1 \\
        0
    \end{pmatrix},
    \qquad
    \vec{x}_2
    =
    \begin{pmatrix}
        3 \\
        0
    \end{pmatrix},
\]

\[
    \vec{\mu}
    =
    \begin{pmatrix}
        2 \\
        0
    \end{pmatrix}.
\]

Berechne die innerhalb-Cluster-Quadratsumme für diesen Cluster.

\begin{examplebox}
Für den ersten Punkt:

\[
    \|\vec{x}_1-\vec{\mu}\|^2
    =
    \left\|
    \begin{pmatrix}
        1 \\
        0
    \end{pmatrix}
    -
    \begin{pmatrix}
        2 \\
        0
    \end{pmatrix}
    \right\|^2
    =
    (-1)^2+0^2
    =
    1.
\]

Für den zweiten Punkt:

\[
    \|\vec{x}_2-\vec{\mu}\|^2
    =
    \left\|
    \begin{pmatrix}
        3 \\
        0
    \end{pmatrix}
    -
    \begin{pmatrix}
        2 \\
        0
    \end{pmatrix}
    \right\|^2
    =
    1^2+0^2
    =
    1.
\]

Also:

\[
    J
    =
    1+1
    =
    2.
\]
\end{examplebox}

\subsection{Aufgabe 4: Silhouette-Wert}

Für einen Datenpunkt sei:

\[
    a_n=2,
    \qquad
    b_n=5.
\]

Berechne den Silhouette-Wert.

\begin{examplebox}
Der Silhouette-Wert ist:

\[
    s_n
    =
    \frac{b_n-a_n}{\max(a_n,b_n)}.
\]

Einsetzen:

\[
    s_n
    =
    \frac{5-2}{\max(2,5)}
    =
    \frac{3}{5}
    =
    0{,}6.
\]

Der Wert ist positiv und relativ hoch.
Der Punkt passt also eher gut zum eigenen Cluster.
\end{examplebox}

\subsection{Aufgabe 5: Hartes oder weiches Clustering}

Bei einem Verfahren gilt für einen Datenpunkt:

\[
    r_{n1}=0{,}2,
    \qquad
    r_{n2}=0{,}7,
    \qquad
    r_{n3}=0{,}1.
\]

Ist das hartes oder weiches Clustering?

\begin{examplebox}
Da die Zugehörigkeiten Werte zwischen \(0\) und \(1\) annehmen und sich zu

\[
    0{,}2+0{,}7+0{,}1=1
\]

summieren, handelt es sich um weiches Clustering.

Bei hartem Clustering wäre genau ein Wert \(1\) und alle anderen \(0\).
\end{examplebox}

\section{Häufige Fehler}

\begin{examtipbox}
Typische Fehler bei Clustering:

\begin{enumerate}
    \item Clustering mit Klassifikation verwechseln.
    \item Labels erwarten, obwohl Clustering unüberwacht ist.
    \item Feature-Skalierung ignorieren.
    \item \(K\) bei \(k\)-means ohne Prüfung festlegen.
    \item Das globale Optimum von \(k\)-means erwarten.
    \item Nur eine Initialisierung verwenden.
    \item Nicht erkennen, dass \(k\)-means kugelförmige Cluster bevorzugt.
    \item Ausreißer ignorieren.
    \item Interne Metriken wie Silhouette als absolute Wahrheit interpretieren.
    \item Cluster automatisch als fachlich bedeutsam ansehen.
    \item PCA-Visualisierung als vollständige Clusterstruktur missverstehen.
    \item Dichtebasierte und zentroidbasierte Methoden verwechseln.
\end{enumerate}
\end{examtipbox}

\section{Zusammenfassung}

\begin{itemize}
    \item Clustering ist unüberwachtes Lernen.
    \item Ziel ist, Datenpunkte in Gruppen ähnlicher Punkte einzuteilen.
    \item Bei \(k\)-means wird die Anzahl der Cluster \(K\) vorgegeben.
    \item Jeder Cluster wird durch ein Zentrum \(\vec{\mu}_k\) beschrieben.
    \item Der Zuordnungsschritt ist:
    \[
        z_n
        =
        \arg\min_k
        \|\vec{x}_n-\vec{\mu}_k\|^2.
    \]
    \item Die \(k\)-means-Zielfunktion ist:
    \[
        J
        =
        \sum_{n=1}^{N}
        \|\vec{x}_n-\vec{\mu}_{z_n}\|^2.
    \]
    \item Das optimale Zentrum eines Clusters ist der Mittelwert seiner Punkte:
    \[
        \vec{\mu}_k
        =
        \frac{1}{N_k}
        \sum_{n:z_n=k}
        \vec{x}_n.
    \]
    \item \(k\)-means konvergiert, aber nur zu einem lokalen Minimum.
    \item Initialisierung ist wichtig.
    \item \(k\)-means++ verbessert die Initialisierung.
    \item Die Wahl von \(K\) kann mit Elbow-Methode oder Silhouette-Koeffizient unterstützt werden.
    \item Hierarchisches Clustering erzeugt eine Clusterhierarchie.
    \item DBSCAN findet dichtebasierte Cluster und kann Ausreißer markieren.
    \item Clustering-Ergebnisse müssen fachlich interpretiert werden.
    \item Gaussian Mixture Models verallgemeinern \(k\)-means probabilistisch.
\end{itemize}

\section{Typische Prüfungsfragen}

\begin{examtipbox}
Mögliche Prüfungsfragen zu diesem Kapitel:

\begin{enumerate}
    \item Was ist Clustering?
    \item Was ist der Unterschied zwischen Clustering und Klassifikation?
    \item Warum ist Feature-Skalierung bei Clustering wichtig?
    \item Definiere die \(k\)-means-Zielfunktion.
    \item Erkläre den Zuordnungsschritt von \(k\)-means.
    \item Leite den Aktualisierungsschritt für die Zentren her.
    \item Warum ist das Zentrum eines Clusters der Mittelwert seiner Punkte?
    \item Warum konvergiert \(k\)-means?
    \item Warum findet \(k\)-means nicht notwendigerweise das globale Optimum?
    \item Warum hängt \(k\)-means von der Initialisierung ab?
    \item Was ist \(k\)-means++?
    \item Was passiert bei leeren Clustern?
    \item Wie kann man \(K\) wählen?
    \item Was ist die Elbow-Methode?
    \item Was misst der Silhouette-Koeffizient?
    \item Welche Grenzen hat \(k\)-means?
    \item Was ist der Unterschied zwischen hartem und weichem Clustering?
    \item Was ist hierarchisches Clustering?
    \item Was ist ein Dendrogramm?
    \item Was sind Single Linkage, Complete Linkage und Average Linkage?
    \item Was ist DBSCAN?
    \item Was ist ein Kernpunkt bei DBSCAN?
    \item Wie unterscheiden sich \(k\)-means und DBSCAN?
    \item Warum ist Clustering schwer zu evaluieren?
    \item Wie hängen \(k\)-means und Gaussian Mixture Models zusammen?
\end{enumerate}
\end{examtipbox}